\documentclass[a4paper]{article}
\usepackage[margin=2.5cm]{geometry}
\usepackage{exam-rtl}

% ----------------------
% Mode Configuration
% ----------------------
\StudentMode
%\TeacherMode

\ShowMarks
\ShowDifficulty
\ShowType
\ShowCorrection
\ExamID{1234}

% --------------------------------------------------
% Document
% --------------------------------------------------
\begin{document}

\begin{center}
    {\Large\bfseries نموذج امتحان}\par\vspace{1em}
\end{center}

% ==================================================
% Multiple Choice
% ==================================================

\begin{Question}{MC}{سهل}{٤}{اختر الإجابة الصحيحة لكل سؤال.}
\begin{MC}[4] 

\MCQuestion{أي كوكب هو الأقرب إلى الشمس وله أقصر فترة مدارية؟}{1}
\MCOption{عطارد}
\MCOption{الزهرة}
\MCOption{المريخ}
\MCOption{الأرض}

\MCQuestion{أي كوكب يدور بعكس اتجاه الدوران المعتاد لباقي كواكب المجموعة الشمسية؟}{2}
\MCOption{عطارد}
\MCOption{الزهرة}
\MCOption{المريخ}
\MCOption{المشتري}

\MCQuestion{أي كوكب يُعرف بامتلاكه أكبر نظام حلقات؟}{3}
\MCOption{نبتون}
\MCOption{أورانوس}
\MCOption{زحل}
\MCOption{المشتري}

\MCQuestion{أي كوكب هو أسرع كواكب النظام الشمسي دورانًا حول نفسه؟}{4}
\MCOption{أورانوس}
\MCOption{نبتون}
\MCOption{زحل}
\MCOption{المشتري}

\end{MC}
\end{Question}

% ==================================================
% True / False
% ==================================================

\begin{Question}{TF}{سهل}{٤}{أجب بصح أو خطأ}
\begin{TF}
\TFQuestion{كوكب الزهرة لديه يوم أطول من سنته.}{صح}
\TFQuestion{زحل هو الكوكب الوحيد الذي يمتلك حلقات في النظام الشمسي.}{خطأ}
\TFQuestion{المشتري أسرع الكواكب دورانا حول محوره.}{صح}
\TFQuestion{البقعة الحمراء العظيمة في كوكب المشتري هي عاصفة ضخمة.}{صح}
\end{TF}
\end{Question}

\begin{Question}{TF}{سهل}{٤}{أجب بصح أو خطأ}
\begin{TF}[options]% أو \begin{TF}[bracket]
\TFQuestion{كوكب عطارد لديه يوم أطول من سنته.}{خطأ}
\TFQuestion{زحل ليس الكوكب الوحيد الذي يمتلك حلقات في النظام الشمسي.}{صح}
\TFQuestion{كوكب الزهرة أبطأ الكواكب دورانا حول محوره.}{صح}
\TFQuestion{البقعة الداكنة العظيمة في كوكب نبتون هي عاصفة عملاقة.}{صح}
\end{TF}
\end{Question}

% ==================================================
% Ordering
% ==================================================

\begin{Question}{Ordering}{متوسط}{٨}{رتب حسب الترتيب الصحيح}
\begin{Ordering}[5]

\OrderQuestion{رتب هذه الكواكب تصاعديًا حسب قربها من الشمس.}
\OrderOption{عطارد}
\OrderOption{الزهرة}
\OrderOption{الأرض}
\OrderOption{المريخ}
\OrderOption{المشتري}
\OrderAnswer
     
\OrderQuestion{رتب هذه الكواكب تنازليًا حسب حجمها.}
\OrderOption{زحل}
\OrderOption{نبتون}
\OrderOption{أورانوس}
\OrderOption{الأرض}
\OrderOption{الزهرة}
\OrderAnswer

\OrderQuestion{رتب هذه الكواكب تنازليًا حسب سرعة دورانها حول محورها.}
\OrderOption{المشتري}
\OrderOption{زحل}
\OrderOption{أورانوس}
\OrderOption{نبتون}
\OrderOption{الأرض}
\OrderAnswer

\OrderQuestion{رتب هذه الكواكب تصاعديًا حسب مدة دورانها حول الشمس.}
\OrderOption{الزهرة}
\OrderOption{المريخ}
\OrderOption{المشتري}
\OrderOption{أورانوس}
\OrderOption{نبتون}
\OrderAnswer

\end{Ordering}
\end{Question}

% ==================================================
% Fill in the Blank
% ==================================================

\begin{Question}{FillBlank}{سهل}{٧}{املأ الفراغ بما يناسبه من الكلمات}
\begin{FillBlank}[true]
\WordBank{
الأرض, عطارد, الحلقات, حرارة, درب التبانة, مدار, الكويكبات, المشتري, أورانوس, مركز, الغلاف الجوي
}
\FillQuestion{حزام \Blank\ هو منطقة تقع بين مداري المريخ والمشتري.}{الكويكبات}
\FillQuestion{كوكب الزهرة هو الكوكب الأكثر \Blank\ في نظامنا الشمسي لأن غلافه الجوي السميك يحبس الحرارة.}{حرارة}
\FillQuestion{كوكب \Blank\ فريد من نوعه لأنه يدور تمامًا على جنبه.}{أورانوس}
\FillQuestion{يقع النظام الشمسي داخل مجرة \Blank[20mm].}{درب التبانة}
\FillQuestion{تقع الشمس في \Blank\ النظام الشمسي.}{مركز}
\FillQuestion{أصغر كوكب من حيث القطر هو \Blank\ وأكبرها هو \Blank\ .}{عطارد ، المشتري}

\end{FillBlank}
\end{Question}

% ==================================================
% Matching
% ==================================================

\begin{Question}{Matching}{متوسط}{٦}{صل ما بين المصطلح وتعريفه الصحيح.}
\begin{Matching}
\MatchPair{لأنه لا يمتلك غلافًا جويًا عمليًا لحبس الحرارة أو توزيعها حول الكوكب.}{تتفاوت درجات الحرارة على سطح عطارد ليلاً ونهاراً بشكل كبير جداً.}
\MatchPair{لأن سطحه مغطى بأكسيد الحديد، أو الصدأ، الذي يعكس الضوء الأحمر.}{المريخ يظهر بلون مائل إلى الحمرة عند رؤيته من الأرض.}
\MatchPair{لأن جاذبية المشتري القوية منعت الحطام الصخري من التكتّل لتكوين كوكب.}{يوجد حزام مميز من الكويكبات بين مداري المريخ والمشتري.}
\MatchPair{لأن دورانه السريع حول محوره وحرارته الداخلية يصنعان عواصف هائلة ومستمرة.}{ تنتشر العواصف الشديدة على المشتري مثل البقعة الحمراء العظيمة أضخم إعصار في النظام الشمسي.}
\MatchPair{لأن تصادمًا هائلًا مع جسم شبيه بكوكب الأرض في مراحله المبكرة أمال محوره.}{أورانوس يدور تقريبًا أفقيًا على جانبه مقارنةً بالكواكب الأخرى.}
\MatchPair{لأن إشعاع الشمس يبخّر جليده مكوّنا تيارات من الغاز والغبار.}{مذنب هالي يكوّن ذيلاً مضيئًا عند اقترابه من الشمس.}
\end{Matching}
\end{Question}

\begin{Question}{Matching}{متوسط}{٦}{صل ما بين المصطلح وتعريفه الصحيح.}
\begin{Matching}[0.2]
\MatchPair{صخرة فضائية تحترق تماماً في غلاف الأرض.}{الشهاب}
\MatchPair{بقايا صخرية ومعدنية لم تحترق وتصطدم بالأرض.}{النيزك}
\MatchPair{جرم يتكون من جليد وصخور وغبار، وله ذيل مضيء.}{المذنب}
\MatchPair{كتل صخرية تقع معظمها بين المريخ والمشتري.}{الكويكبات}
\MatchPair{جرم مضيء يشع ضوءاً وحرارة نتيجة تفاعلات داخله.}{النجم}
\MatchPair{سحابة ضخمة من الغبار والغاز في الفضاء.}{السديم}
\end{Matching}
\end{Question}

% ==================================================
% Short Answer
% ==================================================

\begin{Question}{ShortAnswer}{سهل}{٤}{أكتب جوابا قصيرا.}
\begin{ShortAnswer}[true]
    \ShortQuestion[0]{أي كوكب يُعرف باسم الكوكب الأزرق؟}{نبتون}
    \ShortQuestion[0]{أي كوكب يُعرف باسم الكوكب الأحمر؟}{المريخ}
    \ShortQuestion[0]{أي كوكب يُعرف باسم الكوكب العملاق؟}{المشتري}
    \ShortQuestion[0]{أي كوكب يُعرف بأنه الكوكب التوأم للأرض؟}{الزهرة}
\end{ShortAnswer}
\end{Question}

% ==================================================
% Essay
% ==================================================

\begin{Question}{Essay}{صعب}{٩}{أجب عن الأسئلة التالية بطريقة علمية ومنطقية.}
\begin{Essay}[true]  % أو \begin{Essay}
    \EssayQuestion{بالرغم من أن كوكب عطارد هو الأقرب إلى الشمس، إلا أن كوكب الزهرة هو الأكثر حرارة في النظام الشمسي. فسّر علميًا سبب ذلك؟}{لأن كوكب الزهرة يمتلك غلافًا جويًا سميكًا وكثيفًا يتكون أساسًا من ثاني أكسيد الكربون، مما يؤدي إلى احتباس الحرارة الشديد ، بينما يفتقر عطارد لغلاف جوي حقيقي يحفظ الحرارة.}
    \EssayQuestion[3]{تخيل أنك رائد فضاء هبطت على سطح كوكب المريخ، وحاولت القفز عاليًا. هل ستكون قفزتك أعلى أم أقل ارتفاعًا مقارنة بقفزتك على الأرض؟ علل إجابتك.}{ستكون القفزة أعلى بكثير على المريخ؛ لأن جاذبية المريخ أضعف من جاذبية الأرض (تعادل حوالي 38\% فقط من جاذبية الأرض)، مما يقلل من القوة التي تشد الجسم لأسفل.}
    \EssayQuestion[4]{لماذا تمتلك الكواكب الخارجية (مثل المشتري وزحل) سنوات طويلة جدًا مقارنة بالكواكب الداخلية (مثل الأرض وعطارد)؟}{لسببين رئيسيين: الأول هو أن الكواكب الخارجية تقع على مسافات بعيدة جدًا من الشمس مما يجعل مداراتها حولها أوسع وأطول بكثير، والثاني هو أن سرعتها المدارية (سرعة حركتها حول الشمس) تصبح أبطأ كلما ابتعدت عن مركز جاذبية الشمس.}
\end{Essay}
\end{Question}

% ==================================================
% Reading Passage
% ==================================================

\begin{Question}{ReadingPassage}{متوسط}{٥}{أقرأ القطعة وأجب عن الأسئلة.}
\begin{ReadingPassage}{مذنب هالي}
يُعد مذنب هالي أشهر المذنبات في نظامنا الشمسي، اشتهر بعودته المنتظمة إلى المنطقة الداخلية منه كل 76 سنة تقريبًا. يتكون المذنب من مزيج من الجليد والغبار والصخور، ويُشبّه أحيانًا بـ "كرة ثلجية قذرة". وعندما يقترب من الشمس، يتحول جليده إلى غاز بفعل الحرارة، مُشكّلةً سحابة مضيئة وذيلًا طويلاً لامعًا يشير دائمًا بعيدًا عن الشمس. وقد رآه البشر بعينهم المجردة عبر التاريخ، وآخر مرة ظهر فيها كانت عام 1986، وسيُمكن رؤيته مجددًا في عام 2061. وقد دوّنت الحضارات القديمة ملاحظاته، معتقدة أنها تنبئ بأحداث مصيرية أو تحولات كبرى.
\end{ReadingPassage}

\begin{ShortAnswer}[true]
    \ShortQuestion[1]{كم مرة يعود مذنب هالي إلى المنطقة الداخلية من النظام الشمسي؟}{كل 76 سنة تقريبًا.}
    \ShortQuestion[1]{مما يتكون المذنب مثل مذنب هالي؟}{من الجليد والغبار والصخور.}
    \ShortQuestion[1]{ما الذي يجعل المذنب يطور ذيلًا مضيئًا؟}{تحول جليده إلى غاز بفعل حرارة الشمس.}
    \ShortQuestion[1]{متى ظهر مذنب هالي آخر مرة، ومتى يُتوقع عودته؟}{ظهر آخر مرة عام 1986، ويُتوقع عودته عام 2061.}
\end{ShortAnswer}
\end{Question}

\end{document}